\begin{corollary}[6.3.5]
\label{IV.2.s6.3.5-fr}
Sean $X$, $Y$ dos preesquemas localmente noetherianos y $f:X\to Y$ un
morfismo plano.
\begin{enumerate}
\item[\textup{(i)}] Si, en un punto $x\in X$, se tiene
$\operatorname{coprof}(\mathcal O_x)\leq n$, entonces, poniendo $y=f(x)$, se
tiene $\operatorname{coprof}(\mathcal O_y)\leq n$ y
$\operatorname{coprof}(\mathcal O_x\otimes_{\mathcal O_y}\boldsymbol{k}(y))\leq n$; en
particular, si $\mathcal O_x$ es un anillo de Cohen--Macaulay, también lo son
$\mathcal O_y$ y $\mathcal O_x\otimes_{\mathcal O_y}\boldsymbol{k}(y)$.
\item[\textup{(ii)}] Supongamos, recíprocamente, que
$\mathcal O_x\otimes_{\mathcal O_y}\boldsymbol{k}(y)$ sea un anillo de Cohen--Macaulay.
Entonces, si $\operatorname{coprof}(\mathcal O_y)\leq n$ \textup{(resp.} si
$\mathcal O_y$ es un anillo de Cohen--Macaulay), se tiene
$\operatorname{coprof}(\mathcal O_x)\leq n$ \textup{(resp.} $\mathcal O_x$ es un
anillo de Cohen--Macaulay).
\end{enumerate}
\end{corollary}

Basta aplicar (6.3.2) con $M=A=\mathcal O_y$ y $N=B=\mathcal O_x$.

\begin{corollary}[6.3.6]
\label{IV.2.s6.3.6-fr}
Sea $A$ un anillo de Cohen--Macaulay. Entonces
$A'=A[T_1,\ldots,T_n]$ (siendo $T_i$ indeterminadas) es un anillo de
Cohen--Macaulay.
\end{corollary}

En efecto, si se pone $Y=\operatorname{Spec}(A)$,
$X=\operatorname{Spec}(A')$ y si $f:X\to Y$ es el morfismo canónico, $f$ es
\emph{plano} (puesto que $A'$ es un $A$-módulo libre (2.1.2)) y, para todo $y\in Y$,
$\boldsymbol{k}(y)[T_1,\ldots,T_n]$ es un anillo regular (\textbf 0, 17.3.7), luego \emph{a
fortiori} de Cohen--Macaulay; basta, pues, aplicar (6.3.5).

\begin{corollary}[6.3.7]
\label{IV.2.s6.3.7-fr}
Todo cociente de un anillo de Cohen--Macaulay es universalmente catenario.
\end{corollary}

Esto resulta de (6.3.6), de (5.6.1) y de (\textbf 0, 16.5.12).
